% !TEX root = ../main.tex

\chapter{Introduction}

\section{Continuous Literature Work}

Researchers rarely encounter a paper only once.
A paper may first appear in a recommendation feed, later receive a revised preprint, become relevant to a new project, and finally serve as evidence in a technical decision.
During this process, the reader must preserve at least three kinds of continuity.
First, document continuity identifies which paper and revision supplied a statement.
Second, evidential continuity connects a generated explanation to the text that supports it.
Third, interest continuity records why a paper was selected and how later behavior should---or should not---change the system's view of the user's research interests.

Retrieval-augmented generation (RAG) provides a useful foundation because it couples a generator to retrieved non-parametric memory~\cite{lewis2020rag}.
Nevertheless, retrieval does not by itself prove that an answer is supported by the retrieved text.
Citation-oriented evaluation therefore distinguishes answer quality from citation correctness and completeness~\cite{gao2023alce}.
The distinction becomes more important for long scientific documents: language models can use evidence unevenly depending on where it appears in a long context~\cite{liu2024lost}, and chunk-based retrieval can trade global document structure for local factual detail~\cite{zhao2024longrag}.

Personalization introduces a second source of uncertainty.
Opening, saving, skipping, or questioning a paper can be useful feedback, but no single event is a complete statement of research intent.
Recent work on long-term recommendation describes silent engagement and disengagement as unpaired and potentially ambiguous signals~\cite{zhang2025unpaired}.
Other work shows that editable natural-language user profiles can make recommendations more scrutable~\cite{ramos2024scrutable}.
For MNEME, these observations imply that adaptive interest memory must be inspectable and correctable rather than represented as an invisible, authoritative embedding.

The central thesis is consequently narrower than the title's term ``research agent'' may suggest.
MNEME is not claimed to conduct autonomous science.
It is designed as a mobile literature loop whose actions are bounded by versioned paper evidence, declared behavioral signals, and visible user controls.
This framing turns a feature-rich application plan into a falsifiable research program.

\section{From Capabilities to Technical Challenges}

The project materials describe five product stages: reading, remembering, planning, engaging, and connecting.
Those stages communicate the user journey, but they do not identify the conditions under which the journey is technically trustworthy.
This thesis reformulates the product stages as four technical challenges.

\subsection{TC1: Revision-Correct Paper Representation}

A paper identifier is not sufficient when the underlying document can change.
If a new revision is fetched while summaries, chunks, embeddings, or citations remain tied to an older revision, the interface may display internally inconsistent evidence.
TC1 therefore requires a lineage from logical paper identity to revision identity and from each revision to every derived artifact.
Updates must either preserve a justified artifact or invalidate and regenerate it.

\subsection{TC2: Evidence-Bounded Assistance}

A generated answer can contain a valid paper identifier and still lack support for one or more claims.
Scientific claim verification research makes the evidence unit explicit by pairing atomic claims with rationale spans~\cite{wadden2020scifact}.
Fine-grained attribution work likewise argues that document-level identifiers make it difficult for users to locate exact support~\cite{huang2024grounded}.
TC2 therefore separates four checks: identifier validity, evidence availability, claim--evidence consistency, and evidence coverage.

\subsection{TC3: Adaptive and Inspectable Interest Memory}

Research interests have multiple time scales.
A paper opened during a short course assignment may not represent a stable research direction, whereas repeated saving, questioning, and explicit topic editing may provide stronger evidence.
TC3 must combine explicit and implicit signals while preserving their different meanings.
The system must also expose why an interest changed and permit correction or deletion.

\subsection{TC4: Bounded Mobile Integration}

Mobile integration creates constraints that are not visible in a server-only benchmark.
Network access may be intermittent, background execution is restricted, notifications can interrupt the user, and a small screen can hide evidence behind navigation.
TC4 asks whether TC1--TC3 can be presented as a coherent workflow while making latency, failure, privacy, and uncertainty visible.

\section{Research Questions}

The challenges lead to four research questions:

\begin{description}[leftmargin=2.3cm,style=nextline]
  \item[RQ1] How can a mobile-oriented literature system preserve revision-correct paper identity and section-level provenance during ingestion?
  \item[RQ2] How can retrieval and generation be constrained so that scientific answers remain inspectable against retrieved paper evidence?
  \item[RQ3] How can explicit interests and bounded implicit behavior update a user model without hiding or overstating what the system has learned?
  \item[RQ4] Can these mechanisms be integrated into a usable mobile research loop within stated latency, reliability, and privacy boundaries?
\end{description}

RQ2 and RQ3 are the primary questions.
RQ1 supplies the evidence substrate on which RQ2 depends, while RQ4 tests whether the three mechanisms survive end-to-end integration.

\section{Contributions and Evaluation Contract}

Table~\ref{tab:traceability} states the paper's controlling traceability contract.
At this draft checkpoint, the word ``contribution'' refers to a specified mechanism and evaluation obligation.
It does not imply that the corresponding effectiveness result has already been measured.

\begin{table}[htbp]
\centering
\caption{Traceability from technical challenge to research question, contribution, and evaluation. The status column reflects the evidence supplied for this draft.}
\label{tab:traceability}
\begin{tabular}{P{1.0cm}P{1.0cm}P{4.9cm}P{4.0cm}P{2.4cm}}
\toprule
TC & RQ & Contribution & Evaluation & Current status \\
\midrule
TC1 & RQ1 & C1: revision-aware identity, lineage, and derived-artifact invalidation & E1: identity, update, parsing, and provenance fixtures & \planned \\
TC2 & RQ2 & C2: section-aware retrieval and claim-level source verification & E2: retrieval, answer, citation, abstention, and failure analysis & \pending \\
TC3 & RQ3 & C3: explicit-plus-implicit multi-timescale interest memory with correction & E3: temporal recommendation, ablation, diversity, and correction tests & \pending \\
TC4 & RQ4 & C4: mobile integration with visible evidence, failure, and control & E4: device/system measurements and task-level usability study & \observed{}; \limited \\
\bottomrule
\end{tabular}
\end{table}

The first contribution is a provenance contract: each derived artifact is associated with a logical paper, a specific revision, a transformation configuration, and a content hash.
The second contribution is an evidence-bounded assistance contract: answers are decomposed into claims, claims point to source spans, and verification distinguishes different failure types.
The third contribution is a bounded adaptation contract: behavioral events have declared semantics and weights, multiple time scales remain separable, and users can inspect or correct the resulting state.
The fourth contribution is an integration and evaluation framework that connects these mechanisms to a mobile research workflow.

\section{Scope and Claim Boundary}

The available materials include a controlling thesis outline, acceptance criteria, an architecture design, a story map, an interactive HTML prototype, user-interface and user-experience (UI/UX) mockups, and a final usability presentation.
They do not include a product source repository, an immutable source revision, E1--E3 result artifacts, or participant-level E4 records.
Accordingly, Chapters 2 and 3 describe the design and prototype using \planned{} or \observed{} language.
Chapter 4 specifies complete evaluations and reports only the limited observations supported by the presentation.
Chapter 5 does not claim that the pending experiments succeeded.

\section{Thesis Structure}

Chapter 2 translates the four challenges into functional, non-functional, ethical, and interaction requirements.
Chapter 3 defines the four design contributions and their cross-cutting invariants.
Chapter 4 separates evidence status from evaluation design and specifies E1--E4.
Chapter 5 answers each research question at the strength permitted by the current evidence.

